\chapter{Modulation}
\label{ch:modulation}

A steady, unmodulated sine wave --- a \textbf{carrier} --- conveys no
information beyond its own presence. \textbf{Modulation} is the process
of varying some property of a carrier (its amplitude, frequency, or
phase) in a way that encodes information, so it can be recovered at a
receiver by the inverse process, \textbf{demodulation}.

\section{Why Modulate at All?}

Two practical facts force the use of modulation rather than transmitting
audio (or data) directly: efficient radiation requires an antenna
comparable in size to a wavelength (Chapter~\ref{ch:antennas}), and audio
frequencies (tens of Hz to a few kHz) would need impossibly large
antennas; and multiple simultaneous users need to occupy different parts
of the spectrum without interfering, which is only possible if each
signal is shifted to its own radio-frequency ``slot.''

\section{Amplitude Modulation (AM)}

In \textbf{amplitude modulation}, the carrier's amplitude is varied in
proportion to the modulating (information) signal, while its frequency
stays fixed:
\[
v(t) = \left[V_c + V_m\sin(\omega_m t)\right]\sin(\omega_c t)
\]
The ratio $m = V_m/V_c$ is the \textbf{modulation index} (or
percentage, $\times100\%$); $m=1$ (100\%) is the maximum before the
envelope is driven to zero and distortion (\textbf{overmodulation})
begins.

Expanding the AM equation trigonometrically reveals that AM actually
produces \emph{three} components: the original carrier, plus two new
frequencies --- an \textbf{upper sideband (USB)} at $f_c+f_m$ and a
\textbf{lower sideband (LSB)} at $f_c-f_m$:
\[
v(t) = V_c\sin(\omega_c t) + \frac{mV_c}{2}\cos\!\left[(\omega_c-\omega_m)t\right] - \frac{mV_c}{2}\cos\!\left[(\omega_c+\omega_m)t\right].
\]

\begin{keyidea}[AM occupies twice the information bandwidth]
Both sidebands carry the \emph{same} information (each is a complete,
redundant copy), and the carrier itself carries no information at all
--- yet standard AM transmits both sidebands plus the carrier, using
roughly twice the bandwidth actually needed and wasting most of the
transmitted power in the informationless carrier. This inefficiency is
the direct motivation for single-sideband transmission, below.
\end{keyidea}

\section{Single-Sideband (SSB)}

\textbf{Single-sideband} modulation removes the carrier (using a
\textbf{balanced modulator}) and one of the two redundant sidebands
(using a sharp filter or a phasing technique), transmitting only USB or
LSB. This roughly halves the occupied bandwidth compared to full AM and
concentrates all transmitted power into the one sideband that actually
carries information, making SSB dramatically more efficient than AM for
a given transmitter power --- the reason SSB (not AM) is the dominant
analog voice mode on amateur HF bands today. The receiver must supply
its own locally generated carrier (a \textbf{beat frequency oscillator},
BFO, or product detector reference) to recover audio, and any error in
that reconstructed carrier's frequency shifts the recovered audio pitch
--- the reason SSB signals sound garbled when a receiver is mistuned even
slightly.

\begin{examtip}
By long-standing convention, most amateur voice operation uses lower
sideband (LSB) below 10\,MHz and upper sideband (USB) at 10\,MHz and
above. This is a convention, not a law of physics --- either sideband
can technically carry a voice signal --- but following it is essential
for interoperability with other stations.
\end{examtip}

\section{Frequency Modulation (FM)}

In \textbf{frequency modulation}, the carrier's \emph{frequency} is
varied in proportion to the modulating signal, while amplitude stays
constant:
\[
v(t) = V_c\sin\!\left[\omega_c t + \beta\sin(\omega_m t)\right]
\]
where $\beta$ is the \textbf{modulation index}, related to
\textbf{deviation} $\Delta f$ (the peak frequency shift from the carrier)
by $\beta = \Delta f / f_m$. Because the information rides on frequency
rather than amplitude, FM is largely immune to amplitude-varying noise
and interference (most natural and man-made electrical noise is
amplitude noise) --- the \textbf{FM capture effect} and superior
noise-rejection are why FM is preferred for local VHF/UHF voice
repeater and simplex operation, at the cost of requiring more bandwidth
than AM or SSB for a given audio quality (\textbf{Carson's rule}
estimates occupied bandwidth as $B \approx 2(\Delta f + f_m)$).

\section{Phase Modulation (PM)}

\textbf{Phase modulation} varies the carrier's instantaneous phase in
proportion to the modulating signal. PM and FM are closely related ---
mathematically, FM is what results from phase-modulating a carrier with
the \emph{integral} of the audio signal, and PM is what results from
frequency-modulating a carrier with the \emph{derivative} of the audio
signal --- and in practice many ``FM'' transmitters actually generate
their signal using phase modulation of a crystal-referenced oscillator,
which is easier to keep frequency-stable than direct FM of a
free-running oscillator.

\section{Digital Modulation Concepts}

Digital modes (developed further in Chapter~\ref{ch:digital-modes})
apply the same three basic parameters --- amplitude, frequency, phase ---
to a carrier, but switch between a finite set of discrete states rather
than varying continuously:

\begin{itemize}
\item \textbf{On-off keying (OOK)}: the simplest digital amplitude
modulation, the carrier is simply switched on and off --- the basis of
traditional CW (Morse code) transmission.
\item \textbf{Frequency-shift keying (FSK)}: the carrier is switched
between two (or more) discrete frequencies representing binary states
--- used by RTTY and many amateur digital modes.
\item \textbf{Phase-shift keying (PSK)}: the carrier's phase is switched
between discrete states (e.g., $0\degree$ and $180\degree$ for binary
PSK) --- used by PSK31 and as a building block within more complex
modes.
\item \textbf{Quadrature amplitude modulation (QAM)}: combines discrete
amplitude \emph{and} phase states to encode multiple bits per symbol,
packing more data into the same bandwidth at the cost of requiring
higher signal-to-noise ratio to distinguish the closely spaced states
--- used in modern high-data-rate digital modes.
\end{itemize}

\section{Demodulation}

Recovering the original information requires the inverse operation:

\begin{itemize}
\item \textbf{Envelope/diode detection}: rectifying and low-pass
filtering an AM signal recovers its envelope (the original audio) ---
extremely simple, the basis of the classic ``crystal set.''
\item \textbf{Product detection}: mixing the received signal with a
locally generated reference at (or near) the carrier frequency, then
low-pass filtering, recovers audio from SSB and CW signals; this is the
receive-side counterpart of the balanced modulator used to generate SSB.
\item \textbf{FM discrimination}: circuits (classically, a
frequency-to-voltage converter such as a ratio detector or quadrature
detector; in modern receivers, often a digital phase-locked loop or DSP
algorithm) that convert instantaneous frequency deviation directly into
an output voltage proportional to the original modulating signal.
\end{itemize}

\begin{quickreview}
\item Modulation shifts information onto a radio-frequency carrier,
enabling practical antenna sizes and shared spectrum use.
\item AM varies amplitude, producing a carrier plus two redundant
sidebands; SSB removes the carrier and one sideband for much higher
power/bandwidth efficiency.
\item FM varies frequency (deviation $\Delta f$, index $\beta =
\Delta f/f_m$); it trades bandwidth for strong noise immunity via the
capture effect.
\item PM varies phase and is closely mathematically related to FM.
\item Digital modes apply amplitude/frequency/phase switching between
discrete states (OOK, FSK, PSK, QAM); demodulation reverses the original
modulation process.
\end{quickreview}

\begin{practiceproblems}
\item An AM signal has $V_c = 10\,\text{V}$ and $V_m = 4\,\text{V}$.
Find the modulation index and percentage.
\item Explain why SSB uses roughly half the bandwidth of full-carrier
double-sideband AM for the same audio content.
\item An FM signal has a peak deviation of 5\,kHz and a maximum audio
frequency of 3\,kHz. Estimate its occupied bandwidth using Carson's
rule.
\item Explain, conceptually, why a mistuned SSB receiver produces
garbled, ``Donald Duck''-like audio rather than just a change in pitch
alone.
\item Name the digital modulation scheme used by classic on/off-keyed
Morse code, and explain why it is considered the simplest possible
digital modulation.
\end{practiceproblems}
